\section{社交网络演化史：从内容平台到 AI Native 社交}

Tristan 的经历从音频内容平台、陌生人社交、恋爱游戏、AI 陪伴，到 Elys。这条路径可以看作社交产品的一条演化线：先是内容分发，后来是人和人的匹配，再到虚拟角色提供情绪价值，最后到 AI 分身主动帮真人连接。

\subsection{演化路径}

\noindent\begin{longtable}{p{0.22\textwidth}p{0.34\textwidth}p{0.35\textwidth}}
\toprule
阶段 & 核心机制 & 局限 \\
\midrule
内容平台 & 个性化分发音频/内容 & 连接的是人与内容，不是人与人。 \\
陌生人社交 & 低维标签和即时匹配 & 破冰成本高，关系质量不稳。 \\
恋爱游戏/AI 陪伴 & 内容或 AI 提供情绪价值 & 连接对象可能不是真人。 \\
AI 社交网络 & 分身带 context 主动连接真人 & 隐私、信任、真人比例和冷启动困难。 \\
\bottomrule
\end{longtable}

\begin{knowledgebox}{为什么 Elys 不是从零冒出来的}
Elys 继承了内容推荐、陌生人社交、恋爱游戏和 AI 陪伴的经验，但试图把目标从“消费内容/陪伴”推进到“真人连接”。
\end{knowledgebox}

\subsection{本章小结}

AI Native 社交不是凭空出现，而是内容平台、社交匹配和 AI 陪伴三条线的交汇。Elys 的新意在于让分身主动把 context 转成真人连接。

